\chapter[El Plato Vacío]{El Plato Vacío\\ \large Lo que hoy se desprecia en abundancia, mañana se añora en carestía.}

\elegantimage{15_desperdicio_y_escasez/web/assets/hero.jpg}

\lettrine[lines=3, nindent=0.5em, findent=0.2em]{T}{}ener la fortuna de la prosperidad no debe ser jamás un escudo para la indiferencia ni el despilfarro. Cuando tiramos o despreciamos el alimento, ignorando olímpicamente que millones padecen hambre, quebramos nuestra conexión vital con la providencia...

Este acto de absoluta arrogancia rasga los hilos dorados de la supervivencia. La soberbia de malgastar lo que nutre a los cuerpos no cae en saco roto; el universo lo registra como un profundo desprecio por la vida misma.

Ya que el balance divino exige respeto, quien tira el pan hoy será condenado mañana a sudar sangre, polvo y lágrimas en tierras yerras, solo para rogar por las migajas que en el pasado arrojó al lodo.

\section{El Granero Agotado}
\begin{elegantquote}
\textit{Cuentan que un próspero terrateniente organizaba majestuosos banquetes donde se jactaba ordenando tirar a los perros grandes manjares intactos, tan solo para divertir a sus invitados con su opulencia...}

\textit{Aquellos años de ceguera terminaron abruptamente al girar la Rueda. De pronto, se halló a sí mismo en un paramo desolado en una vida posterior, buscando agua entre grietas secas y rogando por un mísero bocado para sobrevivir a una fatiga devoradora.}

\textit{Fue allí, entre el llanto y el polvo, donde comprendió su error: quien humilla a la tierra desperdiciando el fruto de su vientre, la tierra le humillará cerrándole el sustento.}

\end{elegantquote}

\elegantimage{15_desperdicio_y_escasez/web/assets/art.jpg}

\vspace{1em}
\begin{center}
    \textbf{\Large La abundancia jamás perdona a quien desprecia el pan.}
\end{center}
\vspace{1em}

\section{Análisis Kármico y Traducción}
\begin{wrapfigure}{r}{0.5\textwidth}
  \vspace{-1.5em}
  \begin{center}
  \inlineelegantimage[0.45\textwidth]{15_desperdicio_y_escasez/web/assets/pasaje_original.jpg}
  \end{center}
  \vspace{-1.5em}
\end{wrapfigure}
La abundancia jamás debe ser una excusa para la indiferencia y el desperdicio. Cuando ignoramos el hambre del mundo y despilfarramos los recursos como si carecieran de valor sagrado, reescribimos nuestra conexión con el sustento. Esta arrogancia desencadena en un ciclo posterior la pobreza y el sudor agotador solo para conseguir una porción mínima, enseñando la divina lección de que todo lo que la tierra nos da exige respeto absoluto.

